% !TeX root = ../../Thesis.tex
\begin{table}[htbp]
  \centering
  \small
  \renewcommand{\arraystretch}{1.25}
  \setlength{\tabcolsep}{12pt}
  % The S columns come from siunitx and align the numbers on the decimal marker.
  % Header cells must be braced so that siunitx does not try to parse them as numbers.
  \begin{tabular}{r S[table-format=3.1] S[table-format=1.2] S[table-format=3.1]}
    \toprule
    {\textbf{\ThesisText{Compute nodes}{Rechenknoten}}} &
    {\textbf{\ThesisText{Runtime}{Laufzeit}\ (\unit{\second})}} &
    {\textbf{Speedup}} &
    {\textbf{\ThesisText{Efficiency}{Effizienz}\ (\unit{\percent})}} \\
    \midrule
    1 & 120.0 & 1.00 & 100.0 \\
    2 & 66.7  & 1.80 & 90.0  \\
    4 & 37.5  & 3.20 & 80.0  \\
    8 & 23.5  & 5.11 & 63.9  \\
    \bottomrule
  \end{tabular}
  \caption{\ThesisText{Example scaling results for an illustrative workload.}{Beispielhafte Skalierungsergebnisse für eine illustrative Anwendung.}}
  \label{tab:example-results}
\end{table}
